% Notes of the Montreal seminar, 1970 (cote 129, typed part): archive cote 129, site page 74 (PDF page 75).
% typescript. Draft compared: none.
% AI-transcribed by Claude Opus 5.5 (claude-opus-5-5, adaptive thinking, effort high), Message Batches API, batch mtl129_full (msgbatch_01AGnDKHwy6ATCM7ehUa6263), 2026-11-13; not yet checked by eye.
\scanpage{129/74}
\mspage{\hand{\unsure{III-17|III-19}}}

groupe à fibre connexe et ind-unipotent (i.e. $G^{o}/G^{tor}(1) = G(1)^{o}/G(1)^{mult}$
est unipotent). Cette suite de composition canonique est fonctorielle
en $G$, commute aux changements de base et par passage au dual, on véri-%
fie que, par construction même, on a des isomorphismes canoniques
\[ (G/G^{o})^{*} \simeq (G^{*})^{tor} , \quad (G/G^{tor})^{*} \simeq (G^{*})^{o} . \]

7.3. Si la base $S$ est quelconque, on en peut pas en général écrire
un groupe de BT $G$ comme extension d'un groupe de BT ind-étale $G''$
par un groupe de BT à fibres connexes $G'$.

En effet, si on a une telle extension, on sait que l'on obtient
alors une suite exacte
\[ 0 \to G'(1) \to G(1) \to G''(1) \to 0 . \]

Pour tout $s \in S$, considérons alors le rang séparable de
$G(1)_{s}$. On a:

rang sep $(G(1)_{s})$ = card $G(1)(\overline{s})$ = rang sep $(G'(1)_{s}) \times$ rang sep $(G''(1)_{s})$.

Mais le rang séparable de $G'(1)_{s}$ est $1$ car $G'(1)$ est infi-%
nitésimal et le rang séparable de $G''(1)_{s}$ est une fonction localement cons-%
tante de $s \in S$ car $G''(1)$ est étale sur $S$, de sorte que l'on en dé-%
duit que le rang séparable de $G(1)_{s}$ est une fonction localement constante
de $s$.

Or cela n'est pas vrai en général, comme le montre l'exemple de
la variété elliptique modulaire $A$ sur une courbe $S$ en caractéristique

\mspage{\hand{74}}
